\chapter{Greedy-Max}
\section{Procedimento Greedy}

\subsection{Problema}
\begin{center}
    Consideriamo la seguente istanza di Knapsack-01
\end{center}

\begin{minipage}[t]{0.55\textwidth}
    \centering
    \begin{tabular}{ccc}
        \textbf{OGGETTO} & \textbf{VALORE} & \textbf{INGOMBRO} \\
        \midrule
        $x_1$ & 100 & 30 \\
        $x_2$ & 80  & 25 \\
        $x_3$ & 60  & 20 \\
    \end{tabular}
\end{minipage}
\hfill
\begin{minipage}[t]{0.4\textwidth}
    \textbf{CAPACITÀ ZAINO:}
    \vspace{0.5em}
    
    $C = 45$
\end{minipage}

\vspace{2em}

\noindent \textbf{SOLUZIONE:} \quad $S = \{x_2, x_3\}$

La soluzione si trova utilizzando la programmazione Dinamica

\subsection{Procedimento algoritmo Greedy}
\begin{enumerate}
    \item Si effettua una scelta (quella greedy)
    \item Si risolve un sottoproblema
\end{enumerate}

In un algoritmo greedy facciamo la scelta che sembra migliore in un determinato momento e poi risolviamo il problema che deriva da tale scelta.

La scelta di un algoritmo goloso può dipendere dalle precedenti scelte ma non può dipendere dalle scelte future o dalle soluzioni dei sottoproblemi.

\subsection{Soluzione ottenuta attraverso Greedy}
Nel caso presentato:
\begin{enumerate}
    \item Includo $x_1$ nella soluzione
    \item Il sottoproblema da risolvere sarà $\{x_2, x_3\}$ con capacità residua $= 15$
    \item La soluzione del sottoproblema è $\emptyset$
\end{enumerate}

Applicando la tecnica greedy si ottiene l'insieme $\{x_1\} \cup \emptyset$ \uline{che però non è la soluzione ottimale}.


La tecnica greedy non porta sempre a trovare la soluzione ottima, vale la pena capire quando è percorribile perchè permette di ottenere algoritmi semplici e di facile implementazione

\section{Confronto  tra Programmazione Dinamica e Tecnica Greedy}

\begin{minipage}[t]{0.47\textwidth}
  % Linea orizzontale superiore sinistra
  \rule{\linewidth}{1.2pt}\\[0.8em]
  
  \centering
  {PROGRAMMAZIONE DINAMICA}\\[1.4em]
  
  \raggedright
  Le decisioni vengono prese in base alla risoluzione dei sottoproblemi.\\[1.4em]
  
  \begin{enumerate}
    \item Si risolvono i sottoproblemi
    \item Si effettua una scelta
  \end{enumerate}

  \vspace{1.8em}
  \begin{center}
    \textbf{BOTTOM-UP}
  \end{center}
\end{minipage}%
\hfill
% Linea verticale divisoria
\vrule width 1.2pt
\hfill
\begin{minipage}[t]{0.47\textwidth}
  % Linea orizzontale superiore destra
  \rule{\linewidth}{1.2pt}\\[0.8em]
  
  \centering
  PROGRAMMAZIONE GREEDY\\[1.4em]
  
  \raggedright
  Viene presa la decisione che al momento appare la migliore e poi si risolve un sottopb.\\[1.4em]
  
  \begin{enumerate}
    \item Si effettua una scelta
    \item Si risolve \underline{\textbf{UN}} sottoproblema
  \end{enumerate}

  \vspace{1.8em}
    \begin{center}
        \textbf{SOLITAMENTE TOP-DOWN}    
    \end{center}
\end{minipage}

In entrambe le tecniche vale la proprietà della sottostruttura ottima-

Per capire quando la tecnica Greedy funziona occorre introdurre alcuni concetti.

\section{Sistema Indipendenza (S.I)}

\subsection{Definizione}

Un sistema di indipendenza è una coppia $<E,F>$ dove:

\begin{itemize}
    \item $E$ è un insieme finito ($|E| < \infty$)
    \item $F$ è una famiglia di sottoinsiemi di $E\ (F \subseteq \mathcal{P}(E))$ tali che:
    {
        \large
        $$
        \forall A \in F, \forall B \subseteq A \text{ vale che } B \in F
        $$
    }
\end{itemize}

\subsection{Esempio 1}
$$\langle E,F \rangle$$

\begin{itemize}
    \item $E$ è un insieme finito qualunque
    \item $F = \mathcal{P}(E)$
\end{itemize}

\begin{center}
    È un sistema di indipendenza
\end{center}

\subsection{Esempio 2}
$$
    \langle E,F \rangle
$$

\begin{itemize}
    \item $E \subseteq V$ con $V$ spazio vettoriale e $|E| < \infty$
    \item $F = \{A \subseteq E : A \text{ è un insieme di vettori linearmente indipendenti}\}$
\end{itemize}

\begin{center}
    È un sistema di indipendenza
\end{center}

\newpage
\subsection{Esempio 3}
$$
G = (V,E) \text{un grafo non orientato}
$$

\begin{itemize}
    \item $\langle E,F \rangle$ con $F = \{A \subseteq E : (V,A) \text{ è una foresta}\}$
\end{itemize}

\begin{center}
    È un sistema di indipendenza
\end{center}

Infatti, scelto $A \subseteq F$ (quindi $(V,A)$ è una foresta) e scelto $B \subseteq A$ si ha che $(V,B)$ è una foresta, per cui $B \in F$


\section{Insieme Massimale (Non estendibile)}

\subsection{Definizione}

$A \in F$ è massimale se e solo se $A \in E \backslash A$ vale che $A \cup \{e\} \not \in F$

\section{Problema di massimo associato ad un S.I}

\subsection{Istanza}
\begin{itemize}
    \item $\langle E,F \rangle$ è un sistema di indipendenza
    \item $w: E \to R$
\end{itemize}

\subsection{Soluzione}
{
    \large 
    

    $$
    M \in F \text{ massimale tale che }
    w(M) = \max_{\substack{A \in F \text{ con } A \text{ massiale}}}\{w(A)\}
    $$
}

\subsection{Esempio 4}
\begin{itemize}
    \item $E = \{x_1 , \dots, x_m\}$ e $\forall x_i$ ha un ingombro $ing(x_i)$
    \item $F = \{A \subseteq E : ing(A) \leq C\}$
    \item $w : E \to R^+ \quad \forall x_i\ w(x_i)$ è il valore di $x_i$
\end{itemize}

\subsubsection{Nota}

$\langle E,F \rangle$ è un S.I se:
\begin{itemize}
    \item $A \in F,\ ing(A) \leq C$
    \item $B \subseteq A$
    \item $ing(B) \leq ing(a) \leq C \implies B \in F$
\end{itemize}

Inoltre il problema con 
\begin{itemize}
    \item \textbf{Istanza}: $\langle E,F \rangle \text{ e } w$
    \item \textbf{Soluzione}: $M \in F$ massimale tale che $w(M) = \max_{\substack{A \in F \text{ con } A \text{ massiale}}}\{w(A)\}$
\end{itemize}
è il Knapsack-01

\newpage
\section{Algoritmo Greedy Associato ad un S.I $\langle E,F \rangle$ e $w:E \to R^+$}

\subsection{Pseudocodice}
\begin{algorithmic}
    \Procedure{Greedy-Max}{$E,F,w$} \Comment{$E = \{e_1 , \dots, e_n\}$}
        \State $S = \emptyset$
        \State $Q = [1, \dots, m] = (e_1, \dots, e_m)$
        \State $Q = Sort(Q)$ \Comment{Ordinamente decrescente in baso ai pesi}

        \For{$i = 1$ to $n$}
            \If{$S \cup \{Q[i]\} \in F$} \Comment{Costo: $O(n\log(n)+ n \cdot C(n))$}
                \State $S = S \cup \{Q[i]\}$
            \EndIf
        \EndFor

        \State \Return $S$
    \EndProcedure
\end{algorithmic}

\subsection{Cosa restituisce?}
Greedy-Max ritorna sempre un  sottoinsieme massimale

\subsection{Quando Greedy-Max risolve il problema?}
Greedy-Max risolve i problemi associati a dei \textbf{matroidi}, la cui teoria è utile per determinare quando il metodo greedy genera solozuini ottime, essa non tratta tutti i casi ai quali è possibile applicare il metodo greedy ma consente di analizzare molti casi di interesse pratico